\documentclass[11pt]{article}
\usepackage[a4paper,margin=1in]{geometry}
\usepackage{fontspec}
\usepackage[boldfont=false]{xeCJK}
\setCJKmainfont{FandolSong-Regular.otf}
\usepackage{amsmath}
\usepackage{amssymb}
\usepackage{array}
\usepackage{booktabs}
\usepackage{graphicx}
\usepackage{adjustbox}
\usepackage{enumitem}
\setlist[itemize]{topsep=2pt,partopsep=0pt,parsep=0pt,itemsep=2pt,leftmargin=1.4em}
\setlist[enumerate]{topsep=2pt,partopsep=0pt,parsep=0pt,itemsep=2pt,leftmargin=1.6em}
\usepackage{parskip}
\usepackage{fvextra}
\fvset{breaklines=true,breakanywhere=true,fontsize=\small}
\setlength{\parindent}{0pt}

\begin{document}

\begin{center}
  {\LARGE\bfseries International economic issues}\\[0.35em]
  {\large A Level Economics}
\end{center}
\vspace{0.6em}

\section*{Why countries trade}

\textbf{international trade} 国际贸易 lets countries buy and sell with each other. Trade lets each country \textbf{specialise} 专业化 in what it makes best, so the world makes more in total.

\subsection*{Absolute and comparative advantage}

\begin{itemize}
\item a country has an \textbf{absolute advantage} 绝对优势 in a good if it can make \textbf{more} of it with the same resources than another country.

\item a country has a \textbf{comparative advantage} 比较优势 in a good if it can make it at a \textbf{lower opportunity cost} 机会成本 than another country.

\end{itemize}

Comparative advantage is the key idea: a country should make the goods it gives up least to produce, and trade for the rest. Look at this example, where each country uses the same resources:

\begin{center}
\adjustbox{max width=\linewidth}{%
\begin{tabular}{lll}
\toprule
\textbf{Country} & \textbf{Wine (units)} & \textbf{Cloth (units)} \\
\midrule
Country A & 20 & 10 \\
Country B & 5 & 5 \\
\bottomrule
\end{tabular}}
\end{center}

Country A is better at both (absolute advantage in both). But think about opportunity cost:

\begin{itemize}
\item in A, making 1 cloth means giving up 2 wine.

\item in B, making 1 cloth means giving up 1 wine.

\end{itemize}

So B gives up less to make cloth — B has the comparative advantage in cloth. A has it in wine. If A makes only wine and B makes only cloth, and they trade, \textbf{both} can end up with more than before.

\subsection*{Terms of trade}

The \textbf{terms of trade} 贸易条件 measure the price of a country's exports compared with the price of its imports. If export prices rise faster than import prices, the terms of trade "improve" — the country can buy more imports for the same exports.

\subsection*{Limits of the theory}

The gains are real, but the simple theory assumes no transport costs, no trade barriers, and that costs do not change with output. In the real world these do not fully hold, and specialising in one good is risky if its price falls.

\section*{Protectionism}

\textbf{protectionism} 贸易保护主义 means a government protecting home producers from foreign competition. The main tools:

\begin{itemize}
\item a \textbf{tariff} 关税 — a tax on imports, which raises their price.

\item a \textbf{quota} 配额 — a legal limit on the quantity of a good that may be imported.

\item a \textbf{subsidy} 补贴 to home producers, so they can sell more cheaply than foreign rivals.

\item an \textbf{embargo} 禁运 — a complete ban on trade with a country or in a good.

\item \textbf{administrative barriers} 行政壁垒 — extra paperwork and tough standards that make importing hard.

\end{itemize}

\subsection*{For and against}

Arguments \textbf{for} protection include protecting an \textbf{infant industry} 幼稚产业 (a young industry that needs time to grow), saving jobs, and stopping \textbf{dumping} 倾销 (foreign firms selling below cost to destroy local rivals).

Arguments \textbf{against} are strong: protection raises prices for consumers, protects weak firms so they stay inefficient, and can lead to \textbf{retaliation} 报复 — other countries hitting back with their own barriers.

\subsection*{Effects on stakeholders}

A tariff affects each \textbf{stakeholder} 利益相关者 differently: home producers gain, the government gains tariff revenue, but consumers lose (higher prices, less choice) and foreign producers lose sales.

\section*{The balance of payments}

The \textbf{balance of payments} 国际收支 is a record of all money flows between one country and the rest of the world. Its most-studied part is the \textbf{current account} 经常账户, which has four parts:

\begin{itemize}
\item \textbf{trade in goods} 货物贸易 — exports and imports of physical goods.

\item \textbf{trade in services} 服务贸易 — exports and imports of services (tourism, banking).

\item \textbf{primary income} 初次收入 — income from work and investment abroad (wages, profit, interest).

\item \textbf{secondary income} 二次收入 — transfers with nothing given back, like foreign aid.

\end{itemize}

A \textbf{current account deficit} 经常账户赤字 means a country buys more from abroad than it sells. A \textbf{current account surplus} 经常账户盈余 is the opposite. A long deficit can mean rising debt to other countries; a large surplus can mean the country saves a lot but consumes little.

\section*{Exchange rates}

An \textbf{exchange rate} 汇率 is the price of one currency in terms of another. There are three main systems.

\begin{itemize}
\item a \textbf{floating exchange rate} 浮动汇率 is set freely by the demand for and supply of the currency. If more people want a currency (to buy its exports), its value rises.

\item a \textbf{fixed exchange rate} 固定汇率 is held at a set value by the central bank, which buys and sells the currency to keep it there.

\item a \textbf{managed exchange rate} 管理浮动汇率 mostly floats, but the central bank steps in now and then to steady it.

\end{itemize}

\subsection*{Appreciation and depreciation}

Under floating rates:

\begin{itemize}
\item \textbf{appreciation} 升值 is a rise in the currency's value. Exports become dearer for foreigners, and imports become cheaper.

\item \textbf{depreciation} 贬值 is a fall in the currency's value. Exports become cheaper, and imports become dearer.

\end{itemize}

A depreciation can therefore raise export sales and cut imports, helping the current account — but it can also raise inflation, because imports cost more.

\section*{Correcting a current account deficit}

A government can use three kinds of policy to cut a deficit:

\begin{itemize}
\item \textbf{expenditure-reducing} 支出减少 policies lower total demand (higher taxes or interest rates), so people buy fewer imports. But they also slow growth and raise unemployment.

\item \textbf{expenditure-switching} 支出转换 policies move spending away from imports towards home goods (a tariff, or a lower exchange rate). But they can cause retaliation or inflation.

\item \textbf{supply-side} policies raise the \textbf{competitiveness} 竞争力 of home firms (better skills and technology), so exports rise. These work best in the long run but are slow.

\end{itemize}


\end{document}
